\DocumentMetadata{
  pdfversion=2.0,
  pdfstandard=ua-2,
  lang=en-US,
  tagging=on,
  tagging-setup={math/setup=mathml-SE,tabsorder=structure}
}
\documentclass[11pt,letterpaper]{article}
\usepackage[margin=0.68in,includefoot]{geometry}
\usepackage{newcomputermodern}
\usepackage{amsmath,amscd,mathtools}
\usepackage{tikz,adjustbox}
\providecommand{\square}{\mdlgwhtsquare}
\usepackage[hidelinks]{hyperref}
\hypersetup{
  pdftitle={Analysis PhD Exam, May 2003},
  pdfauthor={University of Florida Department of Mathematics},
  pdfsubject={Graduate mathematics examination},
  pdfdisplaydoctitle=true,
  bookmarksopen=true
}
\setcounter{secnumdepth}{0}
\setlength{\parindent}{0pt}
\setlength{\parskip}{0.28em}
\setlength{\emergencystretch}{3em}
\setlength{\leftmargini}{2.15em}
\setlength{\leftmarginii}{2.65em}
\setlength{\leftmarginiii}{3em}
\pagestyle{empty}
\raggedbottom
\allowdisplaybreaks
\NewTaggingSocketPlug{block/list/label}{examproblem}{%
  \tagstructbegin{tag=Lbl}\tagstructend%
  \tagstructbegin{tag=\LBody}%
  \tagstructbegin{tag=\ProblemHeadingTag,title-o={\ProblemHeadingText}}%
  \tagmcbegin{tag=\ProblemHeadingTag,actualtext={\ProblemHeadingText}}%
  #2%
  \tagmcend%
  \tagstructend%
}
\newenvironment{examproblems}[2]{%
  \def\ProblemHeadingTag{#2}%
  \AssignTaggingSocketPlug{block/list/label}{examproblem}%
  \begin{enumerate}%
  \setcounter{enumi}{#1}%
  \renewcommand{\labelenumi}{\textbf{\theenumi.}}%
  \setlength{\topsep}{0.35em}%
  \setlength{\partopsep}{0pt}%
  \setlength{\itemsep}{0.48em}%
  \setlength{\parsep}{0.18em}%
}{\end{enumerate}}
\newenvironment{examsubparts}{%
  \AssignTaggingSocketPlug{block/list/label}{default}%
  \begin{enumerate}%
  \setlength{\topsep}{0.18em}%
  \setlength{\partopsep}{0pt}%
  \setlength{\itemsep}{0.16em}%
  \setlength{\parsep}{0.1em}%
}{\end{enumerate}}
\newenvironment{examinstructions}{%
  \par\begingroup\itshape
}{%
  \par\endgroup\vspace{0.18em}
}
\makeatletter
\renewcommand\subsection{\@startsection{subsection}{2}{0pt}%
  {-1.25ex plus -.25ex minus -.1ex}%
  {0.55ex plus .1ex}%
  {\normalfont\large\bfseries}}
\renewcommand{\@maketitle}{%
  \newpage\null\vskip -1em%
  {\footnotesize\bfseries
    \href{https://www.ufl.edu/}{University of Florida}, \href{https://math.ufl.edu/}{Department of Mathematics}\par}%
  \vskip -0.5em%
  \tagstructbegin{tag=H1,title={Analysis PhD Exam, May 2003}}%
  \tagpdfparaOff%
  \tagmcbegin{tag=H1}%
    {\LARGE\bfseries\href{https://gma.math.ufl.edu/exams/analysis-phd/}{Analysis PhD Exam}, May 2003\par}%
  \tagmcend%
  \tagpdfparaOn%
  \tagstructend%
  \vskip 0.25em%
  \tagmcbegin{artifact}%
    \hrule height 1pt%
  \tagmcend%
  \vskip 1em%
}
\makeatother
\title{Analysis PhD Exam, May 2003}
\author{}
\date{}
\begin{document}
\maketitle
\thispagestyle{empty}
\begin{examinstructions}
Do All Nine.
\end{examinstructions}
\begin{examproblems}{0}{H2}
\def\ProblemHeadingText{Problem 1.}
\item
Suppose \(X,Y\) are topological vector spaces, \(\Lambda:X\to Y\) is linear, and~\(N\) is a closed subspace of~\(X\). Let \(\pi:X\to X/N\) denote the quotient map. Show, if \(\Lambda(x)=0\) for each \(x\in N\), then there is a unique \(f:X/N\to Y\) satisfying \(\Lambda=f\circ\pi\). Show further, \(f\) is linear and~\(\Lambda\) is continuous if and only if~\(f\) is continuous.
\def\ProblemHeadingText{Problem 2.}
\item
Show, if~\(X\) is an infinite dimensional~\(F\)-space (a topological vector space whose topology is induced by a complete invariant metric), then~\(X\) does not have a countable Hamel basis. A Hamel basis is a basis in the sense of linear algebra.
\def\ProblemHeadingText{Problem 3.}
\item
Is every weakly convergent sequence in \(\ell^1\) strongly convergent? Are the weak and strong topologies on \(\ell^1\) the same?
\def\ProblemHeadingText{Problem 4.}
\item
Let~\(\delta\) denote the distribution defined by \(\delta(\phi)=\phi(0)\) for \(\phi\in\mathcal D(\mathbb R)\). Find the support of \(\delta'\). For which \(f\in C^\infty(\mathbb R)\) is \(f\delta'=0\). Is it possible that an \(f\in C^\infty(\mathbb R)\) can vanish on the support of an \(A\in\mathcal D'(\mathbb R)\) and yet \(fA\ne0\)?
\def\ProblemHeadingText{Problem 5.}
\item
Construct a sequence in \(\mathcal D(\mathbb R)\) which converges to~\(0\) in the topology of \(\mathcal S_1\) (tempered distributions), but not in the topology of \(\mathcal D(\mathbb R)\).
\def\ProblemHeadingText{Problem 6.}
\item
Prove, if~\(A\) is a finite dimensional Banach algebra with unit~\(e\), if \(x,y\in A\), and if \(xy=e\), then \(yx=e\).
\def\ProblemHeadingText{Problem 7.}
\item
Show, if~\(X\) is a compact Hausdorff space, then there is a natural one-one correspondence between closed subsets of~\(X\) and closed ideals in \(C(X)\).
\def\ProblemHeadingText{Problem 8.}
\item
Let \(\mathbb T=\{\gamma\in\mathbb C:|\gamma|=1\}\) denote the unit circle. Let~\(U\) denote the bilateral shift, viewed as the operator \(Uf(\gamma)=\gamma f(\gamma)\) on \(L^2(\mathbb T)\) (with respect to arclength measure). Given~\(\omega\) a measurable subset of \(\mathbb T\), let \(P_\omega\) denote the projection onto the subspace of functions which are supported in~\(\omega\). Does \(P_\omega\) commute with~\(U\)? Find the spectral decomposition of~\(U\).
\def\ProblemHeadingText{Problem 9.}
\item
Let~\(A\) denote the commutative Banach algebra of functions of the form
\[
f(x)=\sum_{n\in\mathbb Z}f_n\exp(inx),\quad\text{with }\sum |f_n|<\infty
\]
 under pointwise multiplication and with the norm \(\lVert f\rVert=\sum |f_n|\). Show, if, for each~\(x\), \(f(x)\ne0\), then~\(f\) is invertible in~\(A\). Is~\(A\) a \(C^*\)-algebra under pointwise complex conjugation?
\end{examproblems}
\end{document}
